\documentclass[11pt]{article}
\usepackage{tikz}
\usepackage{graphicx}
\usepackage[margin=1.0in]{geometry}
\usepackage{etoolbox}
\usepackage{amsmath}

\newcommand{\lessonheader}[2]{%
\begin{center} \Large {MA103: Data-Driven Modeling\\Lesson #1: #2}  \end{center}
\hrule
}

\newcommand{\objectiveitem}[1]{\item #1}
\newcommand{\lessonobjectives}[1]{%
\vspace{-0.1in}
\section*{Lesson Objectives}
\vspace{-0.1in}
\begin{enumerate}
\forcsvlist{\objectiveitem}{#1}
\end{enumerate}
\hrule
}

\newcommand{\subheader}[1]{%
\vspace{-0.1in}
\noindent
\section*{#1}
\vspace{-0.1in}
}

\newcommand{\subsubheader}[1]{%
\vspace{-0.1in}
\noindent
\subsubsection*{#1}
\vspace{-0.1in}
}

\begin{document}
\lessonheader{38}{Objective Functions I}

\lessonobjectives{%
  {Construct an objective function and constraints from context and available data (e.g., costs, capacities, resource limits), with correct units.},
  {Formulate a complete LP algebraically.}%
}

\subheader{Warm-Up: Formulating Constraints}

\noindent Record the general form of a linear programming model. Which part represents the objective function? Which part represents the constraints?

\vfill

\noindent Record the general form of a linear programming model \textit{using matrix vector notation}.

\vfill

\noindent [word problem here]

\noindent Formulate the constraints using function notation. 

\vfill

\noindent Formulate the constraints using matrix vector notation.

\vfill

\newpage

\hrule

\subheader{Exercise: Defend your Model}

As we proceed through the exercise, record observations and tools that are relevant to each phase of the modeling process. Consider how the ideas of \textit{selecting}, \textit{using}, and \textit{assessing} a model fit into the modeling framework we use in this class.
\begin{itemize}
\item FORMULATE

\vfill

\item SOLVE

\vfill

\item INTERPRET

\vfill

\item TEST

\vfill

\end{itemize}

\hrule


\subheader{Wrap-Up: Food for Thought}

\noindent Which model is the best? Does the answer depend on the situation the decision-maker is in?

\vfill

\noindent Is modeling a linear process, or is it circular/iterative? Is this true for both linear and non-linear modeling?

\vfill

\noindent Are there models (or tools for assessing models) that we did not cover in class, but that might have helped our decision maker? For two bonus points, bring in a short (3-5 handwritten sentences) discussion of one model or tool not covered in class that we could have applied to today's exercise. Due Monday, prior to the Predictive Modeling Gate.

\end{document}